\documentclass[aspectratio=169]{beamer}
\usepackage[utf8]{inputenc}
\usepackage[russian]{babel}
\usepackage[T2A]{fontenc}
\usepackage{amsmath,amssymb}
\usepackage{tikz}
\usetikzlibrary{positioning,calc}
\usepackage{listings}
\usepackage{xcolor}

% Выбор темы
\usetheme{Madrid}
\usecolortheme{default}

% Настройка цветов
\definecolor{myblue}{RGB}{0,102,204}
\definecolor{mygreen}{RGB}{0,153,76}
\setbeamercolor{structure}{fg=myblue}
\setbeamercolor{title}{fg=white,bg=myblue}
\setbeamercolor{frametitle}{fg=white,bg=myblue}

% Настройка листингов
\lstset{
  language=[LaTeX]TeX,
  basicstyle=\tiny\ttfamily,
  keywordstyle=\color{blue},
  commentstyle=\color{green!60!black},
  stringstyle=\color{red},
  breaklines=true,
  frame=single,
  numbers=left,
  numberstyle=\tiny\color{gray},
  backgroundcolor=\color{gray!5},
  xleftmargin=5pt,
  xrightmargin=5pt
}

% Информация о презентации
\title[Графика в TikZ]{Построение графических объектов в TikZ}
\subtitle{Графы, функции и фракталы}
\author[Викторов Егор]{Выполнил: [Викторов Егор] \\ Группа: [НПМмд02-24]}
\institute[РУДН]{[РУДН]}
\date{\today}

% Логотип (опционально)
% \logo{\includegraphics[height=1cm]{logo.png}}

\begin{document}

% Титульный слайд
\begin{frame}
  \titlepage
\end{frame}

% Содержание
\begin{frame}{Содержание}
  \tableofcontents
\end{frame}

% Раздел 1: Введение
\section{Введение}

\begin{frame}{Цель работы}
  \begin{block}{Основная цель}
    Изучение возможностей пакета TikZ для создания различных графических объектов в системе \LaTeX
  \end{block}
  
  \vspace{0.5cm}
  
  \begin{columns}[T]
    \begin{column}{0.48\textwidth}
      \textbf{Задачи:}
      \begin{itemize}
        \item Построение графов с узлами
        \item Визуализация функций
        \item Создание фракталов
      \end{itemize}
    \end{column}
    
    \begin{column}{0.48\textwidth}
      \textbf{Инструменты:}
      \begin{itemize}
        \item Пакет TikZ/PGF
        \item Рекурсивные алгоритмы
        \item Математические вычисления
      \end{itemize}
    \end{column}
  \end{columns}
\end{frame}

\begin{frame}{Что такое TikZ?}
  \begin{columns}[c]
    \begin{column}{0.55\textwidth}
      \textbf{TikZ} (TikZ ist kein Zeichenprogramm)
      
      \vspace{0.3cm}
      
      \begin{itemize}
        \item Мощный инструмент для векторной графики
        \item Интеграция с \LaTeX
        \item Программируемость
        \item Точность и масштабируемость
      \end{itemize}
      
      \vspace{0.3cm}
      
      \begin{alertblock}{Преимущество}
        Графика создается кодом, что обеспечивает воспроизводимость и точность
      \end{alertblock}
    \end{column}
    
    \begin{column}{0.4\textwidth}
      \begin{tikzpicture}[scale=0.6]
        \draw[fill=myblue!30] (0,0) circle (1.5);
        \draw[fill=mygreen!30] (1.5,0) circle (1.5);
        \node at (0,0) {\LaTeX};
        \node at (1.5,0) {TikZ};
        \draw[->,ultra thick,red] (0.75,-2) -- (0.75,-0.5);
        \node[below] at (0.75,-2) {Графика};
      \end{tikzpicture}
    \end{column}
  \end{columns}
\end{frame}

% Раздел 2: Задание 1
\section{Задание 1: Граф}

\begin{frame}{Задание 1: Построение графа}
  \begin{block}{Постановка задачи}
    Построить граф с 6 узлами, расположенными по окружности, с различными стилями узлов и ребер
  \end{block}
  
  \vspace{0.3cm}
  
  \begin{columns}[T]
    \begin{column}{0.48\textwidth}
      \textbf{Типы узлов:}
      \begin{itemize}
        \item \textcolor{mygreen}{Зеленые заполненные} (A, C, E)
        \item Двойная обводка (B, D, F)
      \end{itemize}
      
      \vspace{0.2cm}
      
      \textbf{Типы ребер:}
      \begin{itemize}
        \item Пунктирные с весами
        \item Кривые с закруглениями
        \item Гладкие кривые
      \end{itemize}
    \end{column}
    
    \begin{column}{0.48\textwidth}
      \textbf{Используемые техники:}
      \begin{itemize}
        \item Полярные координаты
        \item Пользовательские стили
        \item Параметры кривизны
        \item Позиционирование меток
      \end{itemize}
    \end{column}
  \end{columns}
\end{frame}

\begin{frame}[fragile]{Код: Определение стилей}
\begin{lstlisting}
\tikzset{
  ring/.style={
    circle,
    draw=black,
    thin,
    double=black,
    double distance=1.5pt,
    minimum size=9mm,
    inner sep=0pt,
    fill=white
  },
  green/.style={
    circle,
    draw=black,
    very thick,
    minimum size=9mm,
    inner sep=0pt,
    fill=green!60
  }
}
\end{lstlisting}
\end{frame}

\begin{frame}[fragile]{Код: Размещение узлов}
\begin{lstlisting}
% Узлы в полярных координатах (угол:радиус)
\node[green] (A) at (90:3cm)  {A};
\node[ring]  (B) at (150:3cm) {B};
\node[green] (C) at (210:3cm) {C};
\node[ring]  (D) at (270:3cm) {D};
\node[green] (E) at (330:3cm) {E};
\node[ring]  (F) at (30:3cm)  {F};
\end{lstlisting}

\vspace{0.3cm}

\begin{center}
\begin{tikzpicture}[scale=0.4]
  \draw[gray,dashed] (0,0) circle (3cm);
  \foreach \angle/\label in {90/A,150/B,210/C,270/D,330/E,30/F} {
    \node[circle,fill=myblue,inner sep=2pt] at (\angle:3cm) {};
    \node at (\angle:3.5cm) {\label};
  }
  \draw[->,thick] (0,0) -- (90:3cm);
  \node[right] at (45:1.5cm) {$r=3$};
\end{tikzpicture}
\end{center}
\end{frame}

\begin{frame}[fragile]{Код: Рисование ребер}
\begin{lstlisting}
% Пунктирные линии с метками весов
\draw[dotted,blue,line width=1.5pt] (B) -- 
  node[midway,above,fill=white] {6} (F);
\draw[dotted,blue,very thick] (B) -- 
  node[pos=0.5,left,fill=white] {2} (D);

% Голубая кривая с закругленными углами
\draw[cyan!80!black,very thick,rounded corners=6pt]
  (B) to[out=200,in=160] (C) to[out=-20,in=180] (D);

% Красные гладкие кривые
\draw[red!80!black,very thick,smooth]
  (A) to[out=-60,in=90] 
  node[pos=0.5,right,red,fill=white] {$\sqrt{2}$} (D);
\end{lstlisting}
\end{frame}

\begin{frame}{Результат: Граф}
\begin{center}
\begin{tikzpicture}[>=stealth, every node/.style={font=\small},scale=0.85]
  \tikzset{
    ring/.style={circle,draw=black,thin,double=black,double distance=1.5pt,minimum size=9mm,inner sep=0pt,fill=white},
    green/.style={circle,draw=black,very thick,minimum size=9mm,inner sep=0pt,fill=green!60},
  }
  \node[green] (A) at (90:3cm) {A};
  \node[ring]  (B) at (150:3cm) {B};
  \node[green] (C) at (210:3cm) {C};
  \node[ring]  (D) at (270:3cm) {D};
  \node[green] (E) at (330:3cm) {E};
  \node[ring]  (F) at (30:3cm)  {F};
  \draw[dotted,blue,line width=1.5pt] (B) -- node[midway,above,fill=white] {6} (F);
  \draw[dotted,blue,very thick] (B) -- node[pos=0.5,left,fill=white] {2} (D);
  \draw[dotted,blue,very thick] (F) -- node[pos=0.5,right,fill=white] {4} (D);
  \draw[cyan!80!black,very thick,rounded corners=6pt]
    (B) to[out=200,in=160] (C) to[out=-20,in=180] (D);
  \draw[red!80!black,very thick,smooth]
    (A) to[out=-60,in=90] node[pos=0.5,right,red,fill=white] {$\sqrt{2}$} (D);
  \draw[red!80!black,very thick,smooth]
    (A) to[out=-30,in=90] (E);
  \draw[red!80!black,very thick,smooth]
    (A) to[out=180,in=90] (B);
\end{tikzpicture}
\end{center}
\end{frame}

% Раздел 3: Задание 2
\section{Задание 2: Функции}

\begin{frame}{Задание 2: Графики функций}
  \begin{block}{Постановка задачи}
    Построить графики взаимно обратных функций: $y = e^x$ и $y = \ln(x)$
  \end{block}
  
  \vspace{0.4cm}
  
  \begin{columns}[c]
    \begin{column}{0.48\textwidth}
      \textbf{Экспоненциальная функция:}
      \[y = e^x\]
      \begin{itemize}
        \item Область определения: $(-\infty, +\infty)$
        \item Область значений: $(0, +\infty)$
        \item Проходит через $(0, 1)$
      \end{itemize}
    \end{column}
    
    \begin{column}{0.48\textwidth}
      \textbf{Натуральный логарифм:}
      \[y = \ln(x)\]
      \begin{itemize}
        \item Область определения: $(0, +\infty)$
        \item Область значений: $(-\infty, +\infty)$
        \item Проходит через $(1, 0)$
      \end{itemize}
    \end{column}
  \end{columns}
  
  \vspace{0.3cm}
  
  \begin{alertblock}{Свойство}
    Графики симметричны относительно прямой $y = x$
  \end{alertblock}
\end{frame}

\begin{frame}[fragile]{Код: Построение графиков}
\begin{lstlisting}
% Координатные оси
\draw[->,gray!60] (0,-2.6) -- (0,3.6) node[above] {$y$};
\draw[->,gray!60] (-1.6,0) -- (3.2,0) node[right] {$x$};

% Вспомогательная линия y=1
\draw[gray!25] (-1.3,1) -- (3.0,1);

% График y = e^x
\draw[domain=-1.2:2.2, smooth, variable=\x,
      blue!80!black, very thick]
  plot ({\x},{exp(\x)});

% График y = ln(x)
\draw[domain=0.11:2.8, smooth, variable=\x,
      black!80, very thick]
  plot ({\x},{ln(\x)});
\end{lstlisting}
\end{frame}

\begin{frame}{Результат: Графики функций}
\begin{center}
\begin{tikzpicture}[x=1.5cm,y=1.5cm,scale=0.9]
  \draw[->,gray!60] (0,-2.6) -- (0,3.6) node[above] {$y$};
  \draw[->,gray!60] (-1.6,0) -- (3.2,0) node[right] {$x$};
  \draw[gray!25] (-1.3,1) -- (3.0,1);
  \draw[gray!60] (-0.05,1) -- (0.05,1);
  \node[gray!50,left] at (-1.1,1) {$y=1$};
  \draw[fill=white,draw=gray!60] (0,0) circle(0.06);
  \draw[gray!60] (1,-0.05) -- (1,0.05);
  \node[gray!40,below=3pt] at (1,0) {$x=1$};
  \draw[domain=-1.2:2.2,smooth,variable=\x,blue!80!black,very thick]
    plot ({\x},{exp(\x)});
  \node[blue!80!black, right] at (2.05,3.05) {$y=e^{x}$};
  \draw[domain=0.11:2.8,smooth,variable=\x,black!80,very thick]
    plot ({\x},{ln(\x)});
  \node[black!80,right] at (2.3,0.85) {$y=\ln(x)$};
\end{tikzpicture}
\end{center}
\end{frame}

% Раздел 4: Задание 3
\section{Задание 3: Фрактал}

\begin{frame}{Задание 3: Ковер Серпинского}
  \begin{columns}[c]
    \begin{column}{0.55\textwidth}
      \begin{block}{Что такое фрактал?}
        Самоподобная геометрическая фигура, каждая часть которой повторяет всю фигуру целиком
      \end{block}
      
      \vspace{0.3cm}
      
      \textbf{Ковер Серпинского:}
      \begin{itemize}
        \item Квадрат делится на 9 частей (3×3)
        \item Центральная часть удаляется
        \item Процесс повторяется рекурсивно
      \end{itemize}
      
      \vspace{0.3cm}
      
      \textbf{Фрактальная размерность:}
      \[D = \frac{\ln 8}{\ln 3} \approx 1.893\]
    \end{column}
    
    \begin{column}{0.4\textwidth}
      \begin{tikzpicture}[scale=0.5]
        % Этап 0
        \draw[gray] (0,3) rectangle (2,5);
        \fill[myblue] (0,3) rectangle (2,5);
        \node[below] at (1,2.7) {Этап 0};
        
        % Этап 1
        \draw[gray] (3,3) rectangle (5,5);
        \foreach \i in {0,1,2} {
          \foreach \j in {0,1,2} {
            \pgfmathtruncatemacro{\skip}{(\i==1 && \j==1) ? 1 : 0}
            \ifnum\skip=0
              \fill[myblue] ({3+\i*0.666},{3+\j*0.666}) rectangle ++(0.666,0.666);
            \fi
          }
        }
        \node[below] at (4,2.7) {Этап 1};
      \end{tikzpicture}
    \end{column}
  \end{columns}
\end{frame}

\begin{frame}[fragile]{Алгоритм построения}
\begin{lstlisting}
\newcommand{\carpet}[4]{%
  \pgfmathtruncatemacro{\level}{#4}%
  \ifnum\level=0
    % Базовый случай: рисуем квадрат
    \fill[black] (#1,#2) rectangle +(#3,#3);
  \else
    % Рекурсивный случай
    \pgfmathsetmacro{\newsize}{#3/3}%
    \pgfmathtruncatemacro{\nextlevel}{\level-1}%
    \foreach \i in {0,1,2} {
      \foreach \j in {0,1,2} {
        % Пропускаем центр (i=1, j=1)
        \pgfmathtruncatemacro{\skip}{(\i==1 && \j==1) ? 1 : 0}%
        \ifnum\skip=0
          \pgfmathsetmacro{\newx}{#1 + \i*\newsize}%
          \pgfmathsetmacro{\newy}{#2 + \j*\newsize}%
          \carpet{\newx}{\newy}{\newsize}{\nextlevel}%
        \fi
      }
    }
  \fi
}
\end{lstlisting}
\end{frame}

\begin{frame}{Результат: Ковер Серпинского}
\begin{center}
\begin{tikzpicture}[scale=0.45]
  \newcommand{\carpet}[4]{%
    \pgfmathtruncatemacro{\level}{#4}%
    \ifnum\level=0
      \fill[black] (#1,#2) rectangle +(#3,#3);
    \else
      \pgfmathsetmacro{\newsize}{#3/3}%
      \pgfmathtruncatemacro{\nextlevel}{\level-1}%
      \foreach \i in {0,1,2} {
        \foreach \j in {0,1,2} {
          \pgfmathtruncatemacro{\skip}{(\i==1 && \j==1) ? 1 : 0}%
          \ifnum\skip=0
            \pgfmathsetmacro{\newx}{#1 + \i*\newsize}%
            \pgfmathsetmacro{\newy}{#2 + \j*\newsize}%
            \carpet{\newx}{\newy}{\newsize}{\nextlevel}%
          \fi
        }
      }
    \fi
  }
  \foreach \stage in {0,1,2,3,4} {
    \begin{scope}[xshift=\stage*4cm]
      \draw[gray, thin] (0,0) rectangle (3,3);
      \carpet{0}{0}{3}{\stage}
      \node[below] at (1.5,-0.3) {\small Этап \stage};
    \end{scope}
  }
\end{tikzpicture}
\end{center}

\vspace{0.3cm}

\begin{center}
\begin{tabular}{|c|c|c|}
\hline
\textbf{Этап} & \textbf{Квадратов} & \textbf{Площадь} \\
\hline
0 & $8^0 = 1$ & $1$ \\
1 & $8^1 = 8$ & $8/9$ \\
2 & $8^2 = 64$ & $(8/9)^2$ \\
3 & $8^3 = 512$ & $(8/9)^3$ \\
4 & $8^4 = 4096$ & $(8/9)^4$ \\
\hline
\end{tabular}
\end{center}
\end{frame}

\begin{frame}{Анализ сложности}
  \begin{columns}[T]
    \begin{column}{0.48\textwidth}
      \textbf{Временная сложность:}
      \[T(n) = O(8^n)\]
      
      \vspace{0.3cm}
      
      На каждом уровне:
      \begin{itemize}
        \item Проверяется 9 позиций
        \item Рисуется 8 квадратов
        \item Глубина рекурсии: $n$
      \end{itemize}
    \end{column}
    
    \begin{column}{0.48\textwidth}
      \textbf{Пространственная сложность:}
      \[S(n) = O(n)\]
      
      \vspace{0.3cm}
      
      Стек вызовов:
      \begin{itemize}
        \item Максимальная глубина: $n$
        \item Каждый уровень хранит параметры
        \item Линейная зависимость
      \end{itemize}
    \end{column}
  \end{columns}
  
  \vspace{0.5cm}
  
  \begin{exampleblock}{Практическое ограничение}
    Для этапа 5 потребуется нарисовать $8^5 = 32\,768$ квадратов, что может занять значительное время компиляции
  \end{exampleblock}
\end{frame}

% Раздел 5: Заключение
\section{Заключение}

\begin{frame}{Освоенные навыки}
  \begin{columns}[T]
    \begin{column}{0.48\textwidth}
      \textbf{Технические навыки:}
      \begin{itemize}
        \item Работа с узлами и стилями
        \item Построение различных типов линий
        \item Визуализация функций
        \item Рекурсивное программирование
        \item Работа с координатами
      \end{itemize}
    \end{column}
    
    \begin{column}{0.48\textwidth}
      \textbf{Теоретические знания:}
      \begin{itemize}
        \item Теория графов
        \item Математический анализ
        \item Фрактальная геометрия
        \item Алгоритмы и сложность
        \item Векторная графика
      \end{itemize}
    \end{column}
  \end{columns}
  
  \vspace{0.5cm}
  
  \begin{block}{Главный вывод}
    TikZ — мощный инструмент для создания высококачественной научной графики, интегрированной с \LaTeX
  \end{block}
\end{frame}

\begin{frame}{Преимущества TikZ}
  \begin{columns}[c]
    \begin{column}{0.55\textwidth}
      \begin{itemize}
        \item[\color{mygreen}\checkmark] \textbf{Точность} — математически точные построения
        \item[\color{mygreen}\checkmark] \textbf{Масштабируемость} — векторная графика
        \item[\color{mygreen}\checkmark] \textbf{Интеграция} — единый стиль с документом
        \item[\color{mygreen}\checkmark] \textbf{Программируемость} — автоматизация
        \item[\color{mygreen}\checkmark] \textbf{Воспроизводимость} — код вместо файлов
      \end{itemize}
    \end{column}
    
    \begin{column}{0.4\textwidth}
      \begin{tikzpicture}[scale=0.6]
        \draw[fill=myblue!20] (0,0) -- (3,0) -- (3,3) -- (0,3) -- cycle;
        \draw[fill=mygreen!20] (1.5,1.5) circle (1.2);
        \node at (1.5,1.5) {\Large TikZ};
        \draw[->,ultra thick,red] (1.5,0.3) -- (1.5,-0.5);
        \node[below,red] at (1.5,-0.5) {Качество};
      \end{tikzpicture}
    \end{column}
  \end{columns}
\end{frame}

\begin{frame}{Области применения}
  \begin{center}
  \begin{tikzpicture}[
    box/.style={rectangle,rounded corners,draw=myblue,fill=myblue!10,thick,minimum width=3cm,minimum height=1cm,align=center},
    arrow/.style={->,>=stealth,thick,myblue}
  ]
    \node[box] (tikz) at (0,0) {\textbf{TikZ}};
    
    \node[box] (graphs) at (-3,2.5) {Графы и\\диаграммы};
    \node[box] (math) at (0,2.5) {Математические\\иллюстрации};
    \node[box] (schemes) at (3,2.5) {Схемы и\\блок-схемы};
    
    \node[box] (plots) at (-3,-2.5) {Графики\\функций};
    \node[box] (geo) at (0,-2.5) {Геометрические\\построения};
    \node[box] (fractals) at (3,-2.5) {Фракталы и\\узоры};
    
    \draw[arrow] (tikz) -- (graphs);
    \draw[arrow] (tikz) -- (math);
    \draw[arrow] (tikz) -- (schemes);
    \draw[arrow] (tikz) -- (plots);
    \draw[arrow] (tikz) -- (geo);
    \draw[arrow] (tikz) -- (fractals);
  \end{tikzpicture}
  \end{center}
\end{frame}

\begin{frame}{Дальнейшее развитие}
  \begin{block}{Возможные направления}
    \begin{itemize}
      \item Изучение библиотек TikZ (circuits, automata, mindmap)
      \item 3D-графика с помощью tikz-3dplot
      \item Анимация с использованием animate
      \item Создание интерактивных диаграмм
      \item Разработка собственных библиотек стилей
    \end{itemize}
  \end{block}
  
  \vspace{0.3cm}
  
  \begin{exampleblock}{Рекомендуемые ресурсы}
    \begin{itemize}
      \item Официальная документация TikZ \& PGF (1300+ страниц)
      \item TeXample.net — галерея примеров
      \item TeX Stack Exchange — сообщество
    \end{itemize}
  \end{exampleblock}
\end{frame}

\begin{frame}[standout]
  \Huge Спасибо за внимание!
  
  \vspace{1cm}
  
  \Large Вопросы?
  
  \vspace{1cm}
  
  \normalsize
  \begin{tikzpicture}
    \draw[myblue,ultra thick] (0,0) -- (2,0) -- (1,1.732) -- cycle;
    \node[myblue] at (1,0.577) {\Large \LaTeX};
  \end{tikzpicture}
\end{frame}

% Дополнительные слайды (не входят в основную презентацию)
\appendix

\begin{frame}[fragile]{Приложение: Полный код графа}
\begin{lstlisting}[basicstyle=\fontsize{6}{7}\ttfamily]
\documentclass[border=1cm]{standalone}
\usepackage{tikz}
\usetikzlibrary{positioning,calc}
\begin{document}
\begin{tikzpicture}[>=stealth, every node/.style={font=\small}]
  \tikzset{
    ring/.style={circle,draw=black,thin,double=black,
                 double distance=1.5pt,minimum size=9mm,
                 inner sep=0pt,fill=white},
    green/.style={circle,draw=black,very thick,
                  minimum size=9mm,inner sep=0pt,fill=green!60},
  }
  \node[green] (A) at (90:3cm) {A};
  \node[ring]  (B) at (150:3cm) {B};
  \node[green] (C) at (210:3cm) {C};
  \node[ring]  (D) at (270:3cm) {D};
  \node[green] (E) at (330:3cm) {E};
  \node[ring]  (F) at (30:3cm)  {F};
  \draw[dotted,blue,line width=1.5pt] (B) -- 
    node[midway,above,fill=white] {6} (F);
  \draw[dotted,blue,very thick] (B) -- 
    node[pos=0.5,left,fill=white] {2} (D);
  \draw[dotted,blue,very thick] (F) -- 
    node[pos=0.5,right,fill=white] {4} (D);
  \draw[cyan!80!black,very thick,rounded corners=6pt]
    (B) to[out=200,in=160] (C) to[out=-20,in=180] (D);
  \draw[red!80!black,very thick,smooth]
    (A) to[out=-60,in=90] 
    node[pos=0.5,right,red,fill=white] {$\sqrt{2}$} (D);
  \draw[red!80!black,very thick,smooth]
    (A) to[out=-30,in=90] (E);
  \draw[red!80!black,very thick,smooth]
    (A) to[out=180,in=90] (B);
\end{tikzpicture}
\end{document}
\end{lstlisting}
\end{frame}

\begin{frame}{Приложение: Литература}
  \begin{thebibliography}{99}
    \bibitem{tikz}
    Till Tantau.
    \newblock \textit{The TikZ and PGF Packages}.
    \newblock Manual for version 3.1.10, 2023.
    
    \bibitem{mandelbrot}
    Benoit Mandelbrot.
    \newblock \textit{The Fractal Geometry of Nature}.
    \newblock W. H. Freeman and Company, 1982.
    
    \bibitem{knuth}
    Donald E. Knuth.
    \newblock \textit{The \TeX book}.
    \newblock Addison-Wesley, 1986.
    
    \bibitem{latex}
    Leslie Lamport.
    \newblock \textit{\LaTeX: A Document Preparation System}.
    \newblock Addison-Wesley, 1994.
  \end{thebibliography}
\end{frame}

\end{document}
